\section*{Capítulo \ref{ch:b2:plant-plasticity} --- Adaptações das plantas e plasticidade fenotípica}

\begin{solution}{exo:b2:plant-plasticity:1}
Plasticidade: um \omterm{def:b2:meiosis-heredity:vocabulary}{genótipo}, vários \omterm{def:b2:meiosis-heredity:vocabulary}{fenótipos} conforme o ambiente (\omterm{prop:b2:plant-plasticity:sunshade}{folhas
de sol} e de sombra num mesmo carvalho). \omterm{def:b2:plant-plasticity:plasticity}{Norma de reação}: a curva do
\omterm{def:b2:meiosis-heredity:vocabulary}{fenótipo} de um \omterm{def:b2:meiosis-heredity:vocabulary}{genótipo} em função de uma variável ambiental (espessura
da folha em função da luz). \omterm{def:b2:plant-plasticity:plasticity}{Aclimatação}: um ajuste fisiológico
reversível (o endurecimento do centeio contra a geada no outono).
\omterm{def:b2:plant-plasticity:plasticity}{Adaptação}: um ajuste hereditário selecionado ao longo das gerações (o
metabolismo CAM de um cacto).
\end{solution}

\begin{solution}{exo:b2:plant-plasticity:2}
Coleóptilo intacto iluminado de um lado: curva-se em direção à luz.
Ponta coberta com um capuz opaco: nenhuma curvatura --- a ponta percebe
a luz. Ponta coberta com um capuz transparente: curva-se --- o próprio
capuz é inofensivo. Base protegida por um tubo opaco, ponta exposta:
curva-se --- a percepção está na ponta e a resposta, abaixo dela.
\end{solution}

\begin{solution}{exo:b2:plant-plasticity:3}
\omterm{prop:b2:plant-plasticity:sunshade}{Folha de sol}: pequena, espessa, duas ou três camadas paliçádicas,
cutícula espessa, muitos estômatos, fotossíntese máxima alta e ponto de
compensação alto. \omterm{prop:b2:plant-plasticity:sunshade}{Folha de sombra}: grande, fina, uma camada paliçádica,
mais clorofila por grama, respiração baixa, ponto de compensação
baixo. A escolha é feita enquanto a folha é um primórdio na gema, a
partir da luz que ele recebe; uma folha madura não consegue mudar.
\end{solution}

\begin{solution}{exo:b2:plant-plasticity:4}
As células-guarda absorvem potássio e água, incham e, como suas paredes
internas são mais espessas, se arqueiam e abrem o poro. A luz azul e o
baixo $\mathrm{CO_2}$ interno as abrem pela manhã; o \omterm{def:b2:plant-flowering:dormancy}{ácido abscísico},
vindo das raízes num solo que seca ou da própria folha, as fecha numa
seca, assim como o ar muito seco.
\end{solution}

\begin{solution}{exo:b2:plant-plasticity:5}
$\varphi = R/(R + 0.77\,FR)$ com $FR = 1$: $1.15$: $0.60$; $0.7$:
$0.48$; $0.2$: $0.21$; $0.05$: $0.06$. O limiar $0.4$ fica entre a
clareira ($0.48$) e o dossel ($0.21$): a \omterm{thm:b2:plant-plasticity:phytochrome}{evitação da sombra} é ativada
sob o dossel.
\end{solution}

\begin{solution}{exo:b2:plant-plasticity:6}
$E = 0.25\times 0.025 = \qty{6.25}{mmol.m^{-2}.s^{-1}}$; $A =
0.25\times 100/1.6 = \qty{15.6}{\micro mol.m^{-2}.s^{-1}}$; $A/E =
2.5\times 10^{-3}$ mol de carbono por mol de água, isto é, 400
moléculas de água por carbono, $400\times 18/12 = \qty{600}{g}$ de água
por grama de carbono. Em 12 horas: $6.25\times 10^{-3}\times 43200 =
\qty{270}{mol/m^2}$, \qty{4.9}{L} por metro quadrado.
\end{solution}

\begin{solution}{exo:b2:plant-plasticity:7}
$E = 0.25\times 0.005 = \qty{1.25}{mmol.m^{-2}.s^{-1}}$, $A$
inalterado: $A/E = 0.0125$, 80 moléculas de água por carbono,
\qty{120}{g/g} --- um quinto do custo diurno.
\end{solution}

\begin{solution}{exo:b2:plant-plasticity:8}
$v = 2r^{2}\Delta\rho g/9\eta = 2\times 4\times 10^{-12}\times
400\times 9.81/0.09 = \qty{3.5e-7}{m/s}$: \qty{0.35}{\micro m/s},
cerca de \qty{30}{s} para atravessar a célula. Num clinostato, a
direção da gravidade em relação à célula muda a cada poucos segundos de
sedimentação, de modo que os grãos nunca se acumulam de um lado e o
estímulo médio ao longo do tempo de integração da planta (minutos) é
nulo.
\end{solution}

\begin{solution}{exo:b2:plant-plasticity:9}
Extensão média em duas horas: $0.04$; lados $1.2\times 0.04 = 0.048$
e $0.8\times 0.04 = 0.032$; $\theta = 20\times 0.016/2 = 0.16$
rad, cerca de \qty{9}{\degree}. A curvatura se acumula a $0.08$ rad por
hora; $\pi/2$ leva cerca de \qty{20}{h}.
\end{solution}

\begin{solution}{exo:b2:plant-plasticity:10}
As reservas permitem \qty{50}{cm} de caule. Urtigal: \qty{30}{cm} em
7.5 dias por \qty{30}{mg} --- a plântula alcança a luz com reservas de
sobra. Bosque: \qty{10}{m} estão fora de alcance; a plântula gasta seus
\qty{50}{mg} em meio metro de caule pálido em doze dias e morre
estiolada, ao passo que uma estratégia de tolerância à sombra --- poucas
folhas finas, crescimento lento --- poderia tê-la mantido viva durante
anos.
\end{solution}

\begin{solution}{exo:b2:plant-plasticity:11}
O gelo se forma primeiro nos espaços extracelulares, onde a solução é
mais diluída; a pressão de vapor sobre o gelo é menor que sobre a água
da célula, de modo que a água sai da célula para se juntar ao gelo, a
célula se desidrata e suas membranas colapsam. Uma semana de frio
carrega a célula de açúcares e proteínas protetoras que baixam seu
ponto de congelamento e estabilizam as membranas, e produz proteínas
anticongelantes que limitam o crescimento dos cristais. Manter o
programa ativo o ano todo desvia carbono do crescimento para a proteção
e desacelera a divisão: uma planta permanentemente endurecida é uma
planta pequena.
\end{solution}

\begin{solution}{exo:b2:plant-plasticity:12}
O sinal de vermelho extremo é honesto quando a vizinha pode ser
ultrapassada, e mentiroso quando ela é uma árvore; a semana quente é
honesta em abril e mentirosa em fevereiro; numa cultura densa, as
vizinhas de cada planta são suas iguais, de modo que o sinal é honesto
na forma, mas inútil, e os melhoristas selecionam plantas que o
ignoram. Uma \omterm{def:b2:plant-plasticity:plasticity}{norma de reação} é o registro da frequência com que cada
sinal disse a verdade no passado da linhagem, e ela falha sempre que o
presente difere desse passado.
\end{solution}

\begin{solution}{pb:b2:plant-plasticity:1}
\textbf{1.} Open: $1.15/(1.15 + 0.77) = 0.60$; canopy: $0.2/(0.2 +
0.77) = 0.21$.
\textbf{2.} $\varphi$ baixo: há plantas acima --- alongar-se. Sob uma
tela de sombreamento, $\varphi$ continua $0.60$: um dia escuro, e não
uma vizinha; nenhuma \omterm{thm:b2:plant-plasticity:phytochrome}{evitação da sombra}, apenas um crescimento mais
lento por falta de luz.
\textbf{3.} $1 + 2(0.60 - 0.21) = 1.78$: \qty{2.7}{cm} por dia.
\textbf{4.} \qty{37}{cm} contra \qty{21}{cm}.
\textbf{5.} A curvatura em direção à clareira (fototropismo, por meio do
receptor de luz azul fototropina) e um alongamento mais rápido e a
reorientação das folhas do lado com $R{:}FR$ mais alto (por meio do
fitocromo).
\textbf{6.} As \omterm{prop:b2:plant-plasticity:sunshade}{folhas de sombra} são construídas grandes e finas para
captar a luz escassa a baixo custo; em sol pleno, superaquecem, sofrem
fotoinibição e se queimam, e a planta precisa substituí-las por \omterm{prop:b2:plant-plasticity:sunshade}{folhas
de sol}.
\textbf{7.} \qty{2.7}{cm} por dia numa zona de \qty{15}{mm} são
\qty{0.11}{cm} por hora, uma extensão relativa de $0.074$ por hora;
lado sombreado $1.3\times 0.074 = 0.096$, lado iluminado
$0.7\times 0.074 =
0.052$.
\textbf{8.} $\theta = 15\times(0.096 - 0.052)/2 = 0.33$ rad, cerca de
\qty{19}{\degree} em uma hora.
\textbf{9.} \qty{60}{\degree} são $1.05$ rad: cerca de três horas.
\textbf{10.} $v = 2\times 6.25\times 10^{-12}\times 400\times
9.81/0.09 = \qty{5.5e-7}{m/s}$, \qty{0.55}{\micro m/s}: \qty{22}{s}
para atravessar a célula.
\textbf{11.} O fototropismo, e a própria \omterm{thm:b2:plant-plasticity:phytochrome}{evitação da sombra} aumenta o
ângulo com que o caule se contenta em crescer: a luz é o recurso em
disputa, e um caule que se endireitasse contra a gravidade voltaria
para a sombra.
\textbf{12.} A raiz se curva para baixo até ficar vertical: os
estatólitos sedimentam no lado inferior, a auxina se acumula ali e, numa
raiz, o excesso de auxina inibe o crescimento, de modo que o lado
superior cresce mais depressa. A mesma redistribuição, com a
sensibilidade oposta das células.
\textbf{13.} $E = 0.2\times 0.015 = \qty{3}{mmol.m^{-2}.s^{-1}}$; em
$2\times 10^{-3}$ \unit{m^2} e \qty{50400}{s}: \qty{0.30}{mol},
\qty{5.4}{g} de água por dia.
\textbf{14.} $A = 0.2\times 80/1.6 = \qty{10}{\micro mol.m^{-2}.s^{-1}}$;
por dia, $0.50$ mol de carbono por metro quadrado; $A/E = 3.3\times
10^{-3}$: 300 moléculas de água por carbono, \qty{450}{g/g}.
\textbf{15.} $0.50\times 2\times 10^{-3} = \qty{1}{mmol}$: \qty{12}{mg}
de carbono por dia.
\textbf{16.} Clareira: $E = \qty{6}{mmol.m^{-2}.s^{-1}}$, \qty{10.9}{g}
por dia; eficiência de \qty{900}{g/g}.
\textbf{17.} $E = 0.02\times 0.03 = \qty{0.6}{mmol.m^{-2}.s^{-1}}$
(dez vezes menos); $A = 0.02\times 200/1.6 = \qty{2.5}{\micro
mol.m^{-2}.s^{-1}}$ (quatro vezes menos). A água cai mais: à medida que
a folha reduz seu $\mathrm{CO_2}$ interno, o gradiente para o carbono
aumenta, e o da água não pode aumentar.
\textbf{18.} \qty{4.25}{g} de água, \qty{0.85}{g} a perder; a
\qty{10.9}{g} por dia ($0.78$ g por hora), ela murcha em cerca de uma
hora.
\textbf{19.} A CAM é lenta e exige tecido suculento de reserva; uma
plântula que corre em direção à luz sob um dossel, onde a perda de água
já é pequena, precisa de velocidade, e não de economia.
\textbf{20.} $100/2 = \qty{50}{cm}$.
\textbf{21.} Urtigal: \qty{40}{cm} em 15 dias por \qty{80}{mg} --- ela
alcança a luz. Bosque: \qty{8}{m} são inalcançáveis; ela gasta suas
reservas em meio metro e morre. Melhor no bosque: ficar baixa, formar
folhas finas de sombra e esperar uma clareira --- o modo da faia.
\textbf{22.} Bétula: uma norma íngreme, com forte alongamento a
$\varphi$ baixo, pois as vizinhas de uma pioneira são do seu tamanho.
Faia: uma norma plana, com pouco alongamento e muita tolerância, pois
suas vizinhas são árvores.
\textbf{23.} Desativar a resposta de \omterm{thm:b2:plant-plasticity:phytochrome}{evitação da sombra} --- manter o
fitocromo B ativo ou remover os \omterm{prop:b2:cell-differentiation:transcriptional}{fatores de transcrição} que ele
restringe --- e verificar a época de floração, a germinação e a
resistência do caule, que a mesma via controla.
\textbf{24.} Uma planta que não percebe as vizinhas é ultrapassada
sempre que poderia ter vencido; uma planta que sempre acredita no sinal
se desperdiça sob cada árvore. A norma vencedora aposta no sinal em
proporção à frequência com que, no passado da linhagem, ele esteve
certo.
\textbf{25.} $\varphi = 0.60$ em campo aberto e $0.21$ sob o dossel;
\qty{37}{cm} após 14 dias; \qty{5.4}{g} de água por dia na sombra e
\qty{10.9}{g} na clareira.
\end{solution}
